One important step in rendering participating media is determining attenuation due to scatterimg amd absorbtion through the scene. The total attenuation for photon mapping/splatting methods is usually split into two parts: light-side attenuation and camera-side one. The two are combined at the site where gathering happens (Figure). 

The attenuation from the light side is implicitly stored within the photon's power. Although not calculated in the weight (not explicitly updated in the weight during photon tracing), it is embedded by the free flight sample as its sampling strategy considers the optic thickness. The camera-side attenuation, however, is never calculated before gatherimg happens. Photon mapping and various other methods calculate it through ray-tracing. In our method, we are inspired by methods that handle shadows in rasterization such as opacity shadow map to construct serveral maps for the attenuation value to be retrieved efficiently. 

Opacity Shadow Map (OSM) however is not a method specific for volume rendering but for tackling self-shadowing effect for geometry that consists of abundant small details such as hairs and furs. Like the classic shadow map and the deep shadow map, OSM rasterize the scene from the view point of the light. Unlike the two, it aggregates alpha value instead of depth. OSM cuts the view frustrum from the light into a couple of slices and for each pair of neighboring slices (here we call it slab). It renders all the geometries within the slab without depth test (so it renders everything) to aggregate their alpha value. Overall opacity of any point on the scene (though the method precompute opacity for certain points which they call it shadow sample point) can be then computed with such piecewise step function of alpha value after accumulate the slabs one after another. 

Though not directly applicable, our methodgets inspired by it to construct a pre-computed camera-side attenuation map so that the during shading for the splats (which is determined by this Equation), the camera-side attenuation value can be looked up directly. As a result, the map construction is from the camera's view point, thus deferred rending is taken into account to further speed up the process. And instead of slicing the view frustrum, we use depth peeling to get layers tight to the geometries. Details of how we construct and use such map will be shwon in Chapter 3 in greater details.
